\cotefile{67}
\scanpage{36}
\mspage{10}

\deleted{\unsure{et}}avec les pôles $0,\infty$. D'autre part, l'antipodisme $\sigma'z=-b/\bar z$ \deleted{transforme}\inserted{intervertit} $0$ et $\infty$, et \deleted{\unsure{pour} $z\in C$}\inserted{transforme} $C$ en lui-= car
\[
\begin{aligned}
z\in C&\quad\text{i.e.}\quad z\bar z=b\\
\text{implique}\quad \sigma'z\in C&\quad\text{i.e.}\quad
\underset{\underbrace{(-b/\bar z)(-b/z)}_{b^2/z\bar z}}{\underbrace{(\sigma'z)}\underbrace{\overline{(\sigma'z)}}}=b.
\end{aligned}
\]

Ainsi, $\sigma'$ et $(\sigma,\{0,\infty\})$ définissent la = structure riemannienne sphérique. On est content !

2) $N$ est ext \uline{non centrale} de $\gamma$ par $T$.

Cela signifie que si $\sigma\in N\setminus T$, on aura
\[
\sigma z=a\bar z,\quad\deleted{\unsure{\ne|1}}\quad(a\in\mathbb C^{\unsure{*|\times}}).
\]

\deleted{\unsure{D'où|Oui|Donc}}On en conclut
\[
\sigma^2z=a(\bar a z)=a\bar a z
\]
et la condition $\sigma^2\in T$ signifie
\[
a\bar a=1\quad\text{\deleted{\unsure{d.i.e.|i.e.|d'où}}}\quad a\in\mathbb U.
\]
qui implique aussi $\deleted{\sigma^2=1}$
\[
\sigma^2=1\qquad\text{i.e. l'extension $N$ est \uline{scindée}}
\]
\margin{Ici $N$ est déterminé en termes de $T$, i.e en termes de $\{z,z'\}$, comme $T\cup N^-$, où $N^-$ est formé de \uline{toutes} les anti-involutions de $X$ qui fixent $\{z,z'\}$. Donc il n'est pas question, des \unsure{élts seuls} de \unsure{not.}, d'extraire une métrique riem. sphérique}

Le produit $u_\theta\sigma$ $(\theta\in\mathbb U)$ est
\[
u_\theta\sigma(z)=(\theta a)\bar z\qquad\text{de la \deleted{\unsure{fo}} forme.}
\]
Soit $\alpha\inserted{\in\mathbb U}$ tel que
\[
a=\alpha^2=\alpha/\bar\alpha\qquad\text{donc}
\]
\[
z\in\mathbb C^\sigma\Longleftrightarrow z=\alpha\frac{\bar z}{\bar z}\Longleftrightarrow z/\unsure{a|\alpha}=\overline{(z/\alpha)}\quad\text{i.e.}\quad z\in\alpha\mathbb R
\]
% archive sheet number: 38
